%% kapa: exam=1 round=31 year=2020 period=1 form=A subject=economics subject_ko=경제학원론 seq=8 number=48
%% source: Q-Net 감정평가사 기출문제 (cst00309) artlSeq=5207219 file=extracted/1008/5207219/A형/A형/2020년도 제31회 감정평가사 1차 시험_1교시_A형.pdf
%% conversion: kordoc 4.24.51 -> tools/gen1.py (auto); review=proofread-sonnet; flags=none
\begin{question}{48}
X재 시장의 수요곡선은 $Q_{D}=500-4P$ 이고, 공급곡선은 $Q_{S}=-100+2P$ 이다. 시장 균형에서 정부가 $P=80$의 가격 상한을 설정할 때, (ㄱ)\uline{소비자잉여의 변화}와 (ㄴ)\uline{생산자잉여의 변화}는? (단, $Q_{D}$ 는 수요량, $Q_{S}$ 는 공급량, $P$ 는 가격)
\begin{choices}
  \item ㄱ: 증가, ㄴ: 증가
  \item ㄱ: 증가, ㄴ: 감소
  \item ㄱ: 불변, ㄴ: 불변
  \item ㄱ: 감소, ㄴ: 증가
  \item ㄱ: 감소, ㄴ: 감소
\end{choices}
\end{question}
\answer{2}
